%----------------------------------------------------------------------------
\chapter{Conclusion and future work}
\label{chap:conclusion}
%----------------------------------------------------------------------------

%----------------------------------------------------------------------------
\section{Conclusions}
\label{sec:conclusion-summary}

This thesis investigated whether a language model can generate the abstract of a paper from its recorded reference neighbourhood when the paper-specific information reaches the model only through a soft prompt. To study this, I built an enriched citation dataset, combining ogbn-arxiv with Semantic Scholar citation records, together with the text embeddings and experimental splits that the task requires. I implemented and trained a pipeline that connects neighbourhood estimation by a GNN, prefix conditioning and abstract generation with Qwen3-14B. Finally, I developed an evaluation workflow that compares configurations of the pipeline, checks the behaviour of the evaluation measures and distinguishes fluent writing from agreement with the target's research content.

The implemented system demonstrates that neighbourhood-derived soft prompts can guide the generation of coherent abstracts containing information about the target paper's research context. Learned neighbourhood estimation improves content agreement over averaging the reference embeddings. The evaluation does not, however, establish that the generated abstracts faithfully reconstruct the exact methods or findings of the target papers.

The resulting dataset, generation pipeline and evaluation workflow provide a working basis for further investigation of text generation from graph neighbourhood representations. The implementation will be published on GitHub: \textcolor{red}{[TODO: repository URL]}.

%----------------------------------------------------------------------------
\section{Future work}
\label{sec:future-work}

A natural next step is to replace the autoregressive reader with a masked diffusion language model such as LLaDA~\cite{nie2025llada}. Instead of generating the abstract from left to right, such a model starts from a fully masked response and constructs it iteratively, predicting the masked tokens using bidirectional context~\cite{nie2025llada}. With the same neighbourhood information and the same evaluation, this would test whether a different generation process improves agreement with the target's research problem and task.
